\section{Protocol: a faithful three-phase instantiation}
\label{sec:protocol}

Two instruments carry distinct claims. The three-phase instantiation below
carries the quantitative results of Sections~\ref{sec:synthetic}
and~\ref{sec:realsystem}; a direct execution of the library's own
relations---source unmodified, with the build configuration and a replay harness
added---carries one qualitative claim, that the published code produces the
false proof we predict (Section~\ref{sec:native-run}). Neither instrument is a
benchmark: we report no sensitivity or specificity attributable to
\mbox{MST-wi}.

\paragraph{Three phases}
(1)~\emph{Crawl}: record, per subject and object, whether the subject can reach
the object through the interface, and freeze that record---the model the
predicates will consult. (2)~\emph{Change authorization out of band}: alter a
(subject, object) entitlement through a channel outside the crawled interface;
the frozen model is not updated, because the architecture has no mechanism to
update it. (3)~\emph{Execute the MR} against the current state while the
predicates consult the frozen model. Phase~2 is what the original evaluation
does not exercise, and it is where the reachability stand-in can be wrong.

\paragraph{Adjudicator variants and internal controls}
A raw alarm becomes a verdict through an exception policy, and different policies
are different methods: \textbf{v1} (raw alarms proven); \textbf{v2}
(legitimate-sharing exceptions filtered); \textbf{v3} (state-calibrated---consults
a readable authorization introspection surface, abstains when none is readable);
\texttt{degrade\_all} (the floor: abstain on every alarm);
\texttt{direct\_only} (authorization surface only, access outcome ignored---a
deliberately unfair control); and \texttt{direct\_observe} (a succeeding plain
request checked against the same readable surface). The last consumes the same
observable quantity as the MR family and decides the negative result.

\paragraph{Scenarios}
Seven scenarios separate the reasons an MR can misfire: no shared object (S1,
S2); an in-band legitimate share visible in the framework log (S3); an unshared
plain violation (S4); a read-only share plus a write violation (S5); an
out-of-band share invisible to the framework log (S6); an out-of-band share plus
a genuine violation (S7). Crossed with instantiation mode, enforcement
correctness, channel visibility and introspection availability, this yields
\ResSynRecords{} records (\ResSynErrors{} errors) over \ResSynConfigs{}
configurations.

\paragraph{Statistical treatment}
\label{sec:stats}
The \ResSynRecords{} records are not independent observations: they
deterministically expand \ResDesignCells{} design cells crossed with
\ResDesignMrs{} MRs and two deliberately thrown switches. All intervals (Wilson,
or Newcombe for differences; our own implementation in \texttt{compute\_all.py})
are therefore \emph{descriptive}---magnitude and direction only, no formal
hypothesis-test reading---and no conclusion rests on statistical significance.
The replication ceiling is explicit: one constructed system, \ResDesignCells{}
cells, \ResDesignMrs{} MRs, three real systems. Where a headline count has a
small denominator (the anchored rate, \ResFpEthreeHits{}/\ResFpEthreeTotal{}, is
wide by construction) we lead with the larger-denominator count of the same
phenomenon and never quote a small-denominator rate without its interval.
